% Chapter 4, Topic _Linear Algebra_ Jim Hefferon
%  http://joshua.smcvt.edu/linearalgebra
%  2001-Jun-12
\topic{行列式计算的速度}
对于大型矩阵，用行变换
求行列式通常比用置换展开
快得多。
我们通过求出每种方法各执行多少次操作，来把这一说法
精确化。

要比较两个算法的速度，我们考察当输入数据集的规模增长时，
每个算法所用的时间如何增长。
例如，若把输入的规模增大十倍，
所用时间是增长十倍，
还是增长百倍，或增长千倍？
也就是说，
所用时间是正比于数据集的规模，
还是正比于其规模的平方，或正比于其规模的立方，等等？
所用时间正比于平方的算法
比所用时间正比于立方的算法更快。

先考虑
置换展开公式。\index{determinant!permutation expansion}%
\index{permutation expansion}
\begin{equation*}
   \begin{vmat}
      t_{1,1}  &t_{1,2}  &\ldots  &t_{1,n}  \\
      t_{2,1}  &t_{2,2}  &\ldots  &t_{2,n}  \\
               &\vdots                      \\
      t_{n,1}  &t_{n,2}  &\ldots  &t_{n,n}
   \end{vmat}
   =\!\!\sum_{\text{置换\ }\phi}\!\!\!\!
     t_{1,\phi(1)}t_{2,\phi(2)}\cdots t_{n,\phi(n)}
                                 \deter{P_{\phi}}
   % &=
   % \begin{aligned}[t]
   %    &t_{1,\phi_1(1)}\cdot t_{2,\phi_1(2)}\cdots
   %         t_{n,\phi_1(n)}\deter{P_{\phi_1}}       \\  %[.25ex]
   %    &\hbox{}\quad\hbox{}
   %      +t_{1,\phi_2(1)}\cdot t_{2,\phi_2(2)}\cdots
   %         t_{n,\phi_2(n)}\deter{P_{\phi_2}}       \\
   %    &\hbox{}\quad\hbox{}\vdotswithin{+}             \\
   %    &\hbox{}\quad\hbox{}
   %      +t_{1,\phi_k(1)}\cdot t_{2,\phi_k(2)}\cdots
   %         t_{n,\phi_k(n)}\deter{P_{\phi_k}}
   % \end{aligned}
\end{equation*}
不同的 \( n \)-置换共有
$n!=n\cdot(n-1)\cdots 2\cdot 1$ 个，
因此对一个有 $n$~行的矩阵，这个和式有 $n!$ 项
（且每一项内部有 $n$ 次乘法）。
阶乘函数增长很快：
当 $n$ 仅为 $10$ 时，展开式就已有
$10!=3,628,800$ 项。
注意，与阶乘成比例的增长大于与平方成比例的增长  $n!>n^2$
因为 $n!$ 的前
两个因子相乘得 $n\cdot(n-1)$，当 $n$ 很大
时它近似于 $n^2$，而再乘进更多的因子会使
阶乘更大。
类似地，阶乘函数比~$n^3$ 增长更快，
等等。
因此，使用置换
展开公式、从而所执行的操作次数至少与行数的阶乘一样大的算法，
会非常慢。

相比之下，行化简方法所用时间的增长没有这么快。
% Below is a fragment of row-reduction code.
% It is
% in the computer language \textsc{FORTRAN},
% which is widely used for numeric code.
% % (there are no rows below
% % the \lstinline[style=inline]!N!-th
% % so it is not included in the loop).
% \begin{lstlisting}
% DO 10 ROW=1, N-1
%   PIVINV=1.0/A(ROW,ROW)
%   DO 20 I=ROW+1, N
%     DO 30 J=I, N
%       A(I,J)=A(I,J)-PIVINV*A(ROW,J)
%     30 CONTINUE
%   20 CONTINUE
% 10 CONTINUE
% \end{lstlisting}
下面是用计算机语言 Python 写的
行化简脚本。
（\textit{注：}
这里的代码是朴素的；例如，它没有处理
\lstinline[style=inline]!m(p_row, p_row)! 元素为零的情形。
对一个包含全部检验与子情形的完善版本所做的分析会更繁琐，
但给出的速度结果与此大致相同。）
\begin{lstlisting}
import random

def random_matrix(num_rows, num_cols):
    m = []
    for col in range(num_cols):
        new_row = []
        for row in range(num_rows):
            new_row.append(random.uniform(0,100))
        m.append(new_row)
    return m

def gauss_method(m):
    """Perform Gauss's Method on m.  This code is for illustration only
    and should not be used in practice.
      m  list of lists of numbers; each included list is a row
    """
    num_rows, num_cols = len(m), len(m[0])
    for p_row in range(num_rows):
        for row in range(p_row+1, num_rows):
            factor = -m[row][p_row] / float(m[p_row][p_row])
            new_row = []
            for col_num in range(num_cols):
                p_entry, entry = m[p_row][col_num], m[row][col_num]
                new_row.append(entry+factor*p_entry)
            m[row] = new_row
    return m

response = raw_input('number of rows? ')
num_rows = int(response)
m = random_matrix(num_rows, num_rows)
for row in m:
    print row
M = gauss_method(m)
print "-----"
for row in M:
    print row
\end{lstlisting}
除了执行高斯消元法的例程之外，这个程序还有一个例程，
用来生成一个填满随机数的矩阵（这些数
介于 $0$ 与~$100$ 之间，以便下面可读）。
该程序提示用户
输入行数，生成该规模的随机方阵，
并对其做行化简。
\begin{lstlisting}
$ python gauss_method.py
number of rows? 4
[69.48033741746909, 32.393754742132586, 91.35245787350696, 87.04557918402462]
[98.64189032145111, 28.58228108715638, 72.32273998878178, 26.310252241189257]
[85.22896214660841, 39.93894635139987, 4.061683241757219, 70.5925099861901]
[24.06322759315518, 26.699175587284373, 37.398583921673314, 87.42617087562161]
-----
[69.48033741746909, 32.393754742132586, 91.35245787350696, 87.04557918402462]
[0.0, -17.40743803545155, -57.37120602662462, -97.2691774792963]
[0.0, 0.0, -108.66513774392809, -37.31586824349682]
[0.0, 0.0, 0.0, -13.678536859817994]
\end{lstlisting}  %$

在
\lstinline[style=inline]!gauss_method!
例程内部，对每一行 \lstinline[style=inline]!prow!，例程对其下方的各行执行
$\text{\lstinline[style=inline]!factor!}\cdot\rho_{\text{\lstinline[style=inline]!prow!}}+\rho_{\text{\lstinline[style=inline]!row!}}$。
对其下方的每一行，
这都要对该行中的每个元素进行操作。
这是一个三层嵌套的循环。
所以这个程序的运行时间大约是
矩阵行数的立方。
（\textit{评注。}
我们略过了许多问题。
例如，我们可能会担心程序所用的时间
主要消耗在向内存存储以及
从内存取出元素上，而不是消耗在行变换上。
不过，计算模型的建立超出了
我们的范围。）

如果在程序末尾加上这段代码，
\begin{lstlisting}
def do_matrix(num_rows):
    gauss_method(random_matrix(num_rows, num_rows))

import timeit
for num_rows in [10,20,30,40,50,60,70,80,90,100]:
    s = "do_matrix("+str(num_rows)+")"
    t = timeit.timeit(stmt=s, setup="from __main__ import do_matrix",
                      number=100)
    print "num_rows=", num_rows, " seconds=", t
\end{lstlisting}
那么 Python 就会对程序计时。
下面是一次计时测试运行的输出。
\begin{lstlisting}
num_rows= 10  seconds= 0.0162539482117
num_rows= 20  seconds= 0.0808238983154
num_rows= 30  seconds= 0.248152971268
num_rows= 40  seconds= 0.555531978607
num_rows= 50  seconds= 1.05453586578
num_rows= 60  seconds= 1.77881097794
num_rows= 70  seconds= 2.75969099998
num_rows= 80  seconds= 4.10647988319
num_rows= 90  seconds= 5.81125879288
num_rows= 100  seconds= 7.86893582344
\end{lstlisting}
把这些数据画出来，得到一条三次曲线的一部分。
\begin{center}
  \includegraphics{det/bin/detspeed.pdf}
\end{center}

寻找计算行列式最快的算法
是当前的研究课题。
到目前为止，研究者已找到一些算法，
其运行时间介于行数的平方与立方之间。

置换展开与行变换这两种行列式计算
方法所用时间的对比
说明了这样一点：尽管理论上二者给出同样的答案，
但实践中我们要的是性能最好的那一个。


\begin{exercises}
  % \item[\textit{Most of these presume access to a computer.}]
  \item
    为了对典型矩阵的情形有所了解，
    我们可以利用计算机系统生成随机数的能力
    （当然，这些只是伪随机数，因为它们来自
    某个算法，但它们
    能通过多项合理的随机性统计检验）。
    \begin{exparts}
      \partsitem 用随机数填满一个 $\nbyn{5}$ 数组，比如取
        $\clopinterval{0}{1}$ 范围内的数）。
        看它是否是奇异的。
        把该实验重复几次。
        在这个意义下，奇异矩阵是常见还是罕见？
      \partsitem 给你的计算机代数系统计时，求十个随机数的 $\nbyn{10}$ 数组的行列式。
        求出每个数组的平均用时。
        对 $\nbyn{20}$ 数组、
        $\nbyn{30}$ 数组、\ldots $\nbyn{100}$ 数组
        重复前面的事项，
        并与上面给出的数字比较。
        （注意，当一个数组是奇异的时候，我们有时能很快
        判定这一点，
        例如第一行等于第二行时。
        根据你对第一问的回答，你预计
        奇异方程组会在你的平均值中占很大比重吗？）
      \partsitem 画出输入规模对平均用时的图像。
    \end{exparts}
    \begin{answer}
      你的计时结果一部分取决于所用的计算机代数系统，
      一部分取决于你在其上进行计算的计算机的能力。
      但你应得到一条与所示曲线相似的曲线。
%       \begin{exparts}
%         \partsitem Under Octave, \texttt{rank(rand(5))} finds the
%           rank of a $\nbyn{5}$ matrix whose entries are (uniformly
%           distributed) in the interval $[0..1)$.
%           This loop which runs the test $5000$ times
% \begin{lstlisting}
% octave:1> for i=1:5000
% > if rank(rand(5))<5 printf("That's one."); endif
% > endfor
% \end{lstlisting}
%           produces (after a few seconds) returns the prompt, with no output.

%           The Octave script
% \begin{lstlisting}
% function elapsed_time = detspeed (size)
%   a=rand(size);
%   tic();
%   for i=1:10
%      det(a);
%   endfor
%   elapsed_time=toc();
% endfunction
% \end{lstlisting}
%           lead to this session (obviously, your times will vary).
% \begin{lstlisting}
% octave:1> detspeed(5)
% ans = 0.019505
% octave:2> detspeed(15)
% ans = 0.0054691
% octave:3> detspeed(25)
% ans = 0.0097431
% octave:4> detspeed(35)
% ans = 0.017398
% \end{lstlisting}
%           \partsitem Here is the data (rounded a bit), and the graph.
%             \begin{center}
%               \begin{tabular}{l}
%               \begin{tabular}{r|cccccc}
%                  \textit{matrix rows}
%                     &$15$ &$25$ &$35$ &$45$ &$55$ &$65$  \\
%                  \hline
%                  \textit{time per ten}
%                     &$0.003\,4$
%                     &$0.009\,8$
%                     &$0.067\,5$
%                     &$0.028\,5$
%                     &$0.044\,3$
%                     &$0.066\,3$
%               \end{tabular}                                       \\
%               \qquad\begin{tabular}{r|ccc}
%                  \textit{matrix rows}
%                     &$75$ &$85$ &$95$ \\
%                  \hline
%                  \textit{time per ten}
%                     &$0.142\,8$
%                     &$0.228\,2$
%                     &$0.168\,6$
%               \end{tabular}
%               \end{tabular}
%             \end{center}
%           This data is from an average of twenty runs of the above script,
%           because of the possibility that the randomly chosen matrix
%           happens to take an unusually long or short time.
%           Even so, the timing cannot be relied on too heavily; this is
%           just an experiment.
%           \begin{center}
%             \includegraphics[width=0.4\textwidth]{det/mp/ch4.28}
%           \end{center}
%       \end{exparts}
    \end{answer}
  \item
    用上面讨论的两种方法
    手算下列各行列式。
    \begin{exparts*}
      \partsitem $\begin{vmat}[r]
                    2  &1  \\
                    5  &-3
                  \end{vmat}$
      \partsitem $\begin{vmat}[r]
                    3  &1  &1  \\
                   -1  &0  &5  \\
                   -1  &2  &-2
                  \end{vmat}$
      \partsitem $\begin{vmat}[r]
                    2  &1  &0  &0  \\
                    1  &3  &2  &0  \\
                    0  &-1 &-2 &1  \\
                    0  &0  &-2 &1
                  \end{vmat}$
    \end{exparts*}
    对每种方法，数一数每种情形中
    用到的乘法与除法的次数。
    % (Computer multiplications and divisions take
    %  longer than additions and subtractions, so algorithm
    %  designers worry about them more.)
     \begin{answer}
       操作次数取决于我们具体如何执行这些操作。
       \begin{exparts}
         \partsitem 行列式为 $-11$。
           行化简只需一次倍加，
           含两次乘法
           （$-5/2$ 乘 $2$ 加 $5$，以及 $-5/2$ 乘 $1$ 加 $-3$），
           沿对角线取乘积还需一次乘法。
           置换展开需要两次乘法（$2$ 乘
           $-3$ 与 $5$ 乘 $1$）。
         \partsitem 行列式为 $-39$。
           清点操作次数是例行的。
         \partsitem 行列式为 $4$。
       \end{exparts}
     \end{answer}
  \item 计时例程对 \lstinline[style=inline]!do_matrix!
    的使用有一个缺陷。
    该例程做两件事：生成一个随机矩阵，然后
    对它执行 \lstinline[style=inline]!gauss_method!，而
    返回的计时数字是两者合计的结果。
    请写出只对
    \lstinline[style=inline]!gauss_method! 例程计时的代码。
    \begin{answer}
      你应当只构造矩阵一次，然后对它被多次化简
      进行计时。
      注意，如上所写的 \lstinline[style=inline]!gauss_method!
      会改变输入矩阵，因此要小心，不要用
      该例程的输出做循环，否则除第一次循环外，
      其余各次循环都将在一个阶梯形矩阵上进行。
    \end{answer}
  \item
    你能造出一个使你的计算机
    化简用时最长的 $\nbyn{10}$ 数组吗？
    最短的呢？
    \begin{answer}
      单位矩阵通常化简起来用不了多久。
      对于化简起来很慢的矩阵，在 Python 中你可以试试元素
      相当大的矩阵。
      （一些计算机代数系统还具有生成
      特殊矩阵的能力；可以试试那些。）
      % One way to get started is to compare these under Octave:
      % \texttt{det(rand(10));}, versus
      % \texttt{det(hilb(10));}, versus
      % \texttt{det(eye(10));}, versus
      % \texttt{det(zeroes(10));}.
      % You can time them as in \texttt{tic(); det(rand(10)); toc()}.
    \end{answer}
%   \item
%     Write the rest of the FORTRAN program to do a
%     straightforward implementation of calculating determinants via
%     Gauss's Method.
%     (Don't test for a zero leading entry.)
%     Compare the speed of your code to that used in your computer algebra
%     system.
%     \begin{answer}
%       This is a simple one.
% \begin{lstlisting}
% DO 5 ROW=1, N
%    PIVINV=1.0/A(ROW,ROW)
%    DO 10 I=ROW+1, N
%       DO 20 J=I, N
%          A(I,J)=A(I,J)-PIVINV*A(ROW,J)
%       20 CONTINUE
%    10 CONTINUE
% 5 CONTINUE
% \end{lstlisting}
%     \end{answer}
  \item 一些计算机语言规范要求数组
    按“列”存储，也就是说，整个第一列连续存放，
    然后是第二列，等等。
    所给的代码片段利用了这一点吗？
    或者，利用计算机从连续位置取数更快这一事实，
    能否重写代码使之更快？
    \begin{answer}
      没有，上面的代码是按行处理这些数的。
    \end{answer}
\end{exercises}
\endinput
